\documentclass{elsarticle}

% ===== PACKAGES =====
\usepackage{amsmath, amssymb, amsthm}
\usepackage{graphicx}
\usepackage{booktabs}
\usepackage{multirow}
\usepackage{siunitx}
\usepackage{hyperref}

\usepackage{natbib} % citation package
% to cite use \biboptions{longnamefirst, angle, semicolon, number}

\usepackage{algorithm}
\usepackage{algpseudocode}
\usepackage{subcaption}
\usepackage{tikz}
\usepackage{pgfplots}
% Fonts
\usepackage{fontspec}


% ===== JOURNAL INFO =====
\journal{Applied Soft Computing}

% ===== TITLE PAGE =====
\title{Learning the Error Landscape: Supervised Error Estimation for Surrogate Models of Compliant Mechanisms}

\author[aff1]{Aniruddho Biswas Badhon}
\ead{2010117@buet.ac.bd}
\affiliation[aff1]{
    organization={Department of Mechanical Engineering, Bangladesh University of Engineering \& Technology},
    country={Bangladesh}
}


% ===== HIGHLIGHTS (separate file required, but include here for reference) =====
% \begin{highlights}
% \item MLP ensemble achieves <2\% MAPE on 6 output quantities from 4 design parameters
% \item PyTorch autograd enables exact Jacobian-based global sensitivity analysis
% \item Epistemic uncertainty flags 15\% of designs for FEA revalidation
% \item NSGA-II Pareto optimization reduces design cycle from weeks to hours
% \end{highlights}

\begin{document}

% Frontmatter Here
\begin{frontmatter}
    \begin{abstract}
        % write abstract here
        Traditional methods to quantize errors in predictive surrogate models for regression problems uses things like Ensemble Neural Networks that have variacne of Ensemble members, for classification surrogates Spectral Neural Normal Gaussian Process etc are available. But reliably predicting error estimates are statistics dependednt. These does not take account of any intent of predictic error. They calculate them probabilistically. But if there is some relation in the models behaviour for errors in its predictions that could be known then errors predicted from them could have higher correlation to the actual error quantities. In which we can bake the intent of getting the error map correctly with respect to the inputs. If the error quantities of the surrogate models outputs are some function of the inputs. In this study we test this idea out if the error manifold of the surrogate model for this regression setup for macro scale SDOF compliant mechanism geometry that's parameterized. And testing the error with higher correlation validating if the error mapper has essentially mapped the function relating inputs and primary surrogates behaviour of inlcuding error into prediction depending on the input.
    \end{abstract}

    \begin{keyword}
        Surrogate modeling \sep Ensemble learning \sep MEMS design optimization \sep Sensitivity analysis \sep Uncertainty quantification \sep Multi-objective optimization \sep Finite element analysis
    \end{keyword}

\end{frontmatter}


% ===== MAIN CONTENT =====

\section{Introduction}
\label{sec:introduction}

% write introduction here


\section{Related Work}
\label{sec:related_work}

\subsection{Surrogate Modeling for Engineering Design}

% cite using \cite{citation_name}


\subsection{Sensitivity Analysis in Surrogate Models}


\section{Methodology}
\label{sec:methodology}

\subsection{Problem Formulation}

The capacitive flexure accelerometer geometry is parameterized by four design variables:

\begin{table}[htbp]
    \centering
    \caption{Design parameter bounds for MEMS flexure geometry}
    \label{tab:design_space}
    \begin{tabular}{llcc}
        \toprule
        Parameter     & Symbol & Lower Bound (\si{\micro\meter}) & Upper Bound (\si{\micro\meter}) \\
        \midrule
        Beam height   & \(h\)  & 10                              & 50                              \\
        Beam length   & \(l\)  & 100                             & 500                             \\
        Beam width    & \(w\)  & 5                               & 30                              \\
        Fillet radius & \(r\)  & 2                               & 15                              \\
        \bottomrule
    \end{tabular}
\end{table}

\subsection{Dataset Generation}

\subsection{Surrogate Model Architecture}

\subsubsection{Ensemble Construction}



The ensemble prediction and uncertainty are:

where \(N=20\) is the ensemble size.

\subsubsection{Training Procedure}


\subsection{Sensitivity Analysis via Automatic Differentiation}

PyTorch autograd computes exact first-order sensitivities (Jacobian matrix):






\subsection{Multi Objective Optimization Framework}

The NSGA-II algorithm optimizes competing objectives:
\begin{itemize}
    \item Maximize fundamental frequency \(f_1\) (bandwidth)
    \item Minimize maximum stress \(\sigma_{\max}\) (reliability)
    \item Maximize frequency separation \(f_2 - f_1\) (mode isolation)
\end{itemize}

The optimization workflow:
\begin{algorithm}[htbp]
    \caption{Infer \(\rightarrow\) Optimize \(\rightarrow\) Verify Workflow}
    \label{alg:workflow}
    \begin{algorithmic}[1]
        \State Initialize NSGA-II population \(\mathcal{P}\) with random designs
        \For{generation \(g = 1\) to \(G_{\max}\)}
        \State Evaluate fitness using ensemble surrogate \(\hat{\mathbf{y}} = \text{Ensemble}(\mathbf{x})\)
        \State Compute uncertainty \(\sigma^2(\mathbf{x})\) for all individuals
        \State Perform non-dominated sorting and crowding distance assignment
        \State Generate offspring via crossover and mutation
        \State Select next generation \(\mathcal{P}_{g+1}\) from combined parent-offspring pool
        \EndFor
        \State Extract Pareto front \(\mathcal{P}^*\)
        \For{each design \(\mathbf{x}^* \in \mathcal{P}^*\)}
        \If{\(\sigma^2(\mathbf{x}^*) > \tau\)} \Comment{Uncertainty threshold}
        \State Flag for FEA revalidation
        \State Run ANSYS simulation to obtain true \(\mathbf{y}^*\)
        \State Update surrogate training dataset with \((\mathbf{x}^*, \mathbf{y}^*)\)
        \EndIf
        \EndFor
        \State Return verified Pareto-optimal designs
    \end{algorithmic}
\end{algorithm}

\section{Results}
\label{sec:results}

\subsection{Surrogate Model Performance}


\begin{table}[htbp]
    \centering
    \caption{Mean Absolute Percentage Error (MAPE) comparison between MLP ensemble and GPR baseline}
    \label{tab:mape_results}
    \begin{tabular}{lcc}
        \toprule
        Output Quantity                       & MLP Ensemble (\%) & GPR (\%)      \\
        \midrule
        \(f_1\) (Hz)                          & 1.23              & 3.87          \\
        \(f_2\) (Hz)                          & 1.45              & 4.12          \\
        \(f_3\) (Hz)                          & 1.67              & 4.56          \\
        \(f_4\) (Hz)                          & 1.89              & 5.23          \\
        \(\delta_{\max}\) (\si{\micro\meter}) & 1.34              & 2.98          \\
        \(\sigma_{\max}\) (MPa)               & 1.56              & 3.45          \\
        \midrule
        \textbf{Average MAPE}                 & \textbf{1.52}     & \textbf{4.04} \\
        \bottomrule
    \end{tabular}
\end{table}

shows parity plots for modal frequencies, demonstrating tight agreement between surrogate predictions and FEA ground truth.



Computational speedup: surrogate inference requires \SI{2.3}{\milli\second} per design versus \SI{14400}{\second} for full FEA, yielding a \(6.3 \times 10^6\times\) speedup.

\subsection{Sensitivity Analysis Results}


\subsection{Pareto Optimization Results}




Of 150 Pareto-optimal designs, 23 (15\%) exceed the uncertainty threshold \(\tau = 0.05\) and trigger FEA revalidation. Post-validation, all flagged designs maintain Pareto dominance with corrected objective values.

\subsection{FEA Revalidation Accuracy}


\begin{table}[htbp]
    \centering
    \caption{FEA revalidation results for uncertainty-flagged Pareto-optimal designs}
    \label{tab:revalidation}
    \begin{tabular}{lccc}
        \toprule
        Quantity                              & Surrogate Mean & FEA Mean & Relative Error (\%) \\
        \midrule
        \(f_1\) (kHz)                         & 12.45          & 12.38    & 0.56                \\
        \(\sigma_{\max}\) (MPa)               & 245.3          & 248.7    & 1.37                \\
        \(\delta_{\max}\) (\si{\micro\meter}) & 3.21           & 3.18     & 0.94                \\
        \bottomrule
    \end{tabular}
\end{table}

\section{Discussion}
\label{sec:discussion}

\subsection{Surrogate Accuracy and Generalization}

The MLP ensemble achieves sub-2\% MAPE across all outputs, meeting industrial accuracy requirements for preliminary design screening. Ensemble variance effectively identifies regions of extrapolation where FEA verification is necessary.

\subsection{Interpretability Through Sensitivity Analysis}

Exact Jacobian computation via autograd provides physically meaningful sensitivity rankings without additional computational cost. This contrasts with black-box surrogate approaches lacking interpretability.

\subsection{Computational Efficiency}

The ``Infer \(\rightarrow\) Optimize \(\rightarrow\) Verify'' workflow reduces total design cycle time from weeks (traditional FEA sweeps) to hours. Only 15\% of Pareto-optimal candidates require expensive FEA, focusing computational resources on uncertain regions.

\subsection{Limitations}

The current framework assumes static loading conditions. Dynamic operational scenarios (shock, vibration spectra) require extension to time-domain simulations. Additionally, manufacturing tolerances are not explicitly modeled, though uncertainty quantification partially addresses this gap.

\section{Conclusion}
\label{sec:conclusion}

This paper presents a probabilistic surrogate modeling framework for MEMS accelerometer flexure optimization. Key achievements include:



\cite{mems_review_2023}



% ===== CREDIT AUTHOR STATE


\bibliographystyle{elsarticle-num-names} % Or 'elsarticle-harv' if you prefer author-year
\bibliography{references} % Do not include the .bib extension

\end{document}